\documentclass[11pt]{article}
\usepackage{amsmath,amssymb,amsthm,booktabs,longtable,hyperref}
\usepackage[margin=1in]{geometry}
\newtheorem{theorem}{Theorem}
\newtheorem{lemma}{Lemma}
\newcommand{\Rea}{\operatorname{Re}}
\title{A strict strip for Ehrhart roots in dimensions six and seven}
\author{Principia Math In Progress}
\date{29 September 2026}
\begin{document}\maketitle
\begin{abstract}
We prove that every root $z$ of the Ehrhart polynomial of a reflexive
lattice polytope of dimension six or seven satisfies $-3<\Rea z<2$.
This proves the nonreal-root assertion in MathDB 330132, namely
Conjecture 6.4 of the 2015 Heged\H{u}s--Higashitani--Kasprzyk manuscript.
The proof applies to a larger cone of palindromic binomial-basis
polynomials. Its only computation is an exact determinant expansion:
all 112 coefficients used in the proof are printed in the appendix and
verified by an independent integer-arithmetic program.
\end{abstract}

\section*{Task 1: the original problem and the stronger result}
For a six- or seven-dimensional reflexive polytope $P$, the question asks
whether a root $-1/2+u+iv$ of $L_P$, with $u,v\ne0$, necessarily has
$|u|<5/2$. This is Conjecture 6.4 on page 15 of \cite{HHKpre}.
The final article \cite{HHK} omits the preprint's conjecture section;
its page 5 still explicitly leaves the six- and seven-dimensional strip
case unproved. We address the preprint's complete assertion, with the
stronger conclusion that also includes real roots.

\begin{theorem}\label{thm}
Let $d\in\{6,7\}$ and
\[
 F(z)=\sum_{j=0}^d h_j\binom{z+d-j}{d},\qquad h_j=h_{d-j},
\]
where
\[
 h_0>0,\quad h_1\ge h_0,\quad h_2\ge h_1,\quad h_3\ge h_1.
\]
Then every root of $F$ satisfies $-3<\Rea z<2$.
In particular, every root of the Ehrhart polynomial of a reflexive
lattice polytope in either dimension satisfies this strict strip bound.
\end{theorem}

The last assertion follows from the standard Ehrhart expansion and
Hibi's lower bound: the $h^*$-vector of a reflexive polytope is
palindromic, $h_0=h_d=1$, $h_1\ge1$, and $h_j\ge h_1$ for
$2\le j\le d-1$. These facts and their original references are stated
on page 2 of \cite{HHK}. No unimodality assumption is made.

Write
\begin{equation}\label{params}
 h_0=a,\quad h_1=a+b,\quad h_2=a+b+c,\quad h_3=a+b+e,
 \qquad a>0,\quad b,c,e\ge0.
\end{equation}
Define degree-six polynomials
\[
 Q_6(x)=6!F(x+2),\qquad
 Q_7(x)=\frac{7!F(x+2)}{2x+5}.
\]
The quotient in dimension seven is a polynomial: binomial symmetry gives
$F(-1-z)=(-1)^dF(z)$, so $F(-1/2)=0$ when $d=7$.
In both cases write $Q_d(x)=\sum_{r=0}^6 q_r x^{6-r}$.
Expanding the binomial products gives the following exact coefficient
matrices, whose rows are indexed by $r=0,\ldots,6$:
\begin{equation}\label{matrix}
 (q_0,\ldots,q_6)^T=C_d(a,b,c,e)^T,
\end{equation}
\[
 C_6=\begin{pmatrix}
 7&5&2&1\\105&75&30&15\\1015&575&200&85\\5775&2625&750&225\\
 19138&6980&1598&274\\34440&9900&1740&120\\25920&5760&720&0
 \end{pmatrix},\qquad
 C_7=\begin{pmatrix}
 4&3&1&1\\60&45&15&15\\760&423&127&85\\5100&2355&645&225\\
 20836&7542&1744&274\\46680&12960&2220&120\\45360&9072&1008&0
 \end{pmatrix}.
\]
In particular $q_0,q_6>0$ throughout \eqref{params}.

\section*{Excluding crossings of the imaginary axis}
For a real sextic $Q(x)=q_0x^6+q_1x^5+\cdots+q_6$, define
\begin{equation}\label{delta}
 D(Q)=\det\begin{pmatrix}
 q_1&q_3&q_5&0&0\\
 q_0&q_2&q_4&q_6&0\\
 0&q_1&q_3&q_5&0\\
 0&q_0&q_2&q_4&q_6\\
 0&0&q_1&q_3&q_5
 \end{pmatrix}.
\end{equation}
This is the fifth Hurwitz determinant, but the proof below needs no
stability criterion as a black box.

\begin{lemma}\label{axis}
If $q_6\ne0$ and $D(Q)\ne0$, the polynomial $Q$ has no root on the
imaginary axis.
\end{lemma}
\begin{proof}
Zero is not a root. If $iy$ were a root for a nonzero real $y$, then
$t=y^2$ would be a common root of
\[
 E(t)=-q_0t^3+q_2t^2-q_4t+q_6,\qquad
 O(t)=q_1t^2-q_3t+q_5.
\]
Consequently the nonzero vector $(t^4,t^3,t^2,t,1)^T$ would lie in the
kernel of
\[
 K=\begin{pmatrix}
 -q_0&q_2&-q_4&q_6&0\\0&-q_0&q_2&-q_4&q_6\\
 q_1&-q_3&q_5&0&0\\0&q_1&-q_3&q_5&0\\0&0&q_1&-q_3&q_5
 \end{pmatrix}.
\]
Row permutations and column sign changes give $\det K=-D(Q)$,
contradicting its nonzero value.
\end{proof}

For the two matrices \eqref{matrix}, exact expansion of \eqref{delta}
has the form
\begin{equation}\label{certificate}
 D(Q_d)=g_d\sum_{i+j+\ell+m=5} A^{(d)}_{i,j,\ell,m}
              a^i b^j c^\ell e^m,
 \qquad g_6=3888000,\quad g_7=21168000.
\end{equation}
The appendix gives every coefficient. Each is a positive integer;
the smallest is $1176$ for $d=6$ and $216$ for $d=7$.
There are exactly $\binom{8}{3}=56$ monomials in each expansion.
In particular the coefficient of $a^5$ is positive, so
\begin{equation}\label{pos}
 D(Q_6)>0,\qquad D(Q_7)>0
 \quad(a>0,\ b,c,e\ge0).
\end{equation}
This is a finite exact algebraic certificate valid on the entire
unbounded parameter domain, not extrapolation from numerical tests.
It can be verified by substituting \eqref{matrix} into \eqref{delta}
and using the determinant's permutation formula. The accompanying
\texttt{verify.py} does this using only integer arithmetic; it also
reconstructs \eqref{matrix} directly from the binomial products and
checks $\det K=-D(Q_d)$ independently.

\section*{A stable base polynomial and the full argument}
At $b=c=e=0$, all $h_j=a$. The polynomial binomial summation identity gives
\[
 a^{-1}F(z)=\sum_{j=0}^d\binom{z+d-j}{d}
 =\binom{z+d+1}{d+1}-\binom{z}{d+1}.
\]
Every root of this difference has real part $-1/2$. Indeed, put
$z=w-1/2$ and $s_k=k+1/2$, $0\le k\le d$.
The two products to be compared, up to the common factor $(d+1)!$,
are $\prod_k(w+s_k)$ and $\prod_k(w-s_k)$. If $\Rea w>0$, then
$|w+s_k|>|w-s_k|$ for every $k$, so the products cannot be equal.
If $\Rea w<0$, all inequalities reverse. This proves the assertion,
including the possibility that one product is zero.
Thus all roots of $Q_6$ and $Q_7$ at this base point have real part
$-5/2$; removal of the root $z=-1/2$ in the odd case does not affect this.

Fix any parameters \eqref{params} and replace $(b,c,e)$ by
$(tb,tc,te)$, $0\le t\le1$. Along this continuous path, $q_0>0$ and
$q_6>0$, and \eqref{pos} and Lemma~\ref{axis} exclude every imaginary-axis
root. The degrees stay equal to six. After division by $q_0$, the
coefficients vary continuously on a compact interval and remain bounded,
so the roots have a uniform bound. Continuity of polynomial roots
(equivalently, local invariance of the root count by the argument
principle) implies that the number of roots in each open half-plane
cannot change along the path. Initially all six have negative real
part, so the same is true at $t=1$.

Therefore every root of $F$ has real part less than $2$, including the
removed root $-1/2$ when $d=7$. The identity
$F(-1-z)=(-1)^d F(z)$ also sends every root $z$ to $-1-z$.
Applying the upper bound to this reflected root gives $\Rea z>-3$.
This proves Theorem~\ref{thm} and the requested strict inequality
$|\Rea(z)+1/2|<5/2$.\hfill$\square$

\section*{Four proposed follow-up tasks}
\begin{enumerate}
\item[Task 2.] Determine the sharp supremum of
$|\Rea z+1/2|$ among roots with both nonzero imaginary part and
$\Rea z\ne-1/2$, separately for the two coefficient cones in
Theorem~\ref{thm}.
\item[Task 3.] Decide the half-strip assertion in dimensions eight and
nine, first for palindromic coefficient vectors satisfying Hibi's lower
bounds, and then for vectors realized by reflexive lattice polytopes.
\item[Task 4.] Characterize which boundary coefficient directions that
approach the extremal strip in dimensions six and seven can be realized
by reflexive polytopes, and obtain any resulting sharper geometric bound.
\item[Task 5.] Find an analytic factorization or positivity representation
for the determinant in \eqref{delta} that explains the two positive
coefficient expansions and identifies the dimensions in which it persists.
\end{enumerate}

\begin{thebibliography}{9}
\bibitem{HHKpre} G.~Heged\H{u}s, A.~Higashitani, A.~M.~Kasprzyk,
\emph{Ehrhart polynomial roots of reflexive polytopes},
\href{https://arxiv.org/abs/1503.05739v1}{arXiv:1503.05739v1} (2015),
Conjecture~6.4, p.~15.
\bibitem{HHK} G.~Heged\H{u}s, A.~Higashitani, A.~M.~Kasprzyk,
\emph{Ehrhart polynomial roots of reflexive polytopes},
Electronic Journal of Combinatorics 26(1) (2019), P1.38,
\href{https://doi.org/10.37236/7780}{doi:10.37236/7780}.
\end{thebibliography}
\clearpage
\input{certificate-tables.tex}
\end{document}
